%% ============================================================
%%  TRA-SAE — Mixed-domain scientific reasoning with a 4B LLM
%%  Wiley NJD template (USG.cls). Target: Expert Systems (Wiley).
%%  Compile in this folder: pdflatex -> bibtex -> pdflatex x2
%%  Numbers: single-pass (one greedy decoding pass), no verifier-gated retry.
%% ============================================================
\documentclass[APA,twocolumn]{USG}

\usepackage{anyfontsize}
\usepackage{colortbl}
\usepackage{xcolor}
\usepackage{listings}
\lstset{basicstyle=\fontsize{8}{9.5}\selectfont\ttfamily,breaklines=true,
        columns=fullflexible,keepspaces=true,frame=none,aboveskip=4pt,belowskip=4pt}
\graphicspath{{./}{./images/}}

\newcommand{\modelname}{TRA-SAE}
\newcommand{\basemodelname}{Qwen3.5-4B}
\newcommand{\dataset}{EXACT~2026}

\articletype{ORIGINAL ARTICLE}%
%% History dates are inserted by the publisher at production.
%\received{}
%\revised{}
%\accepted{}
\journal{Expert Systems}
\volume{}
\copyyear{2026}
\startpage{1}
\articledoi{}

%% APA (apacite) does not provide natbib-style \citep/\citet; alias them when
%% absent so the same source compiles under the APA style and the Chicago fallback.
\makeatletter
\@ifundefined{citep}{\let\citep\cite}{}
\@ifundefined{citet}{\let\citet\citeA}{}
\makeatother

\begin{document}

\title{\modelname{}: A Low-Cost Retrieval-Augmented Expert Reasoning
       System for Mixed Physics--Logic Question Answering with a 4B
       Language Model}

\author[1]{Anonymous Author(s)}

\authormark{ANONYMOUS \textsc{et al.}}
\titlemark{Expert Reasoning for Physics--Logic QA}

\address[1]{\orgdiv{Affiliation withheld for review}, \orgname{Institution},
\orgaddress{\state{}, \country{}}}

\corres{Corresponding author. \email{withheld@review.org}}

\keywords{large language models | expert reasoning systems | knowledge engineering |
scientific question answering | retrieval-augmented generation | curriculum fine-tuning |
controlled ablation study | explainable educational question answering}

\abstract[ABSTRACT]{Scientific and educational question answering couples
numerical physical reasoning with symbolic logical inference, yet
resource-constrained settings often preclude large proprietary models or
extensive reinforcement learning. This paper presents \modelname{}, a low-cost,
retrieval-augmented expert reasoning system that combines dense retrieval, a
structured response schema, a curriculum-trained 4B open model, and a symbolic
verifier (Z3 and SI-unit checking) in a single physics-and-logic pipeline, using
the \dataset{} challenge as a testbed. Rather than proposing a new algorithm, it
is a controlled attribution of which standard components pay off at 4B scale. The training curriculum combines general
supervised fine-tuning (SFT), a logic-focused stage, and Group Relative Policy
Optimization (GRPO) with lightweight format, correctness, and unit-consistency
rewards, at a training cost of US\$18 (full pipeline below US\$64) on a single
A100-80GB. Every system is evaluated under a single greedy decoding pass. On a
217-sample validation split, overall accuracy rises from 29.95\% to 47.47\%
(63.12\% physics, 18.42\% logic). A controlled decomposition attributes $+9.22$
percentage points (pp) of the gain to retrieval and a further, borderline
$+5.07$~pp to SFT, whereas GRPO reward shaping, a bespoke unit-consistency
reward, self-consistency, and tool augmentation add no statistically separable
benefit. External transfer to MMLU physics and FOLIO logic indicates the
curriculum generalises rather than overfitting the testbed. Despite the modest
absolute accuracy, the study identifies what works---retrieval over weight
training---and what fails: formal logic remains the principal bottleneck at 4B
scale, with most residual errors traced to answer extraction.}

\maketitle

%% ============================================================
\section{Introduction}\label{sec1}
%% ============================================================

Large language models (LLMs) achieve strong results on mathematical and
logical reasoning benchmarks \citep{Wei2022,Wang2023sc,Kojima2022}. Building a
dependable reasoning system under resource constraints, where larger proprietary
models are unavailable, is less a matter of a single new algorithm than of
knowledge engineering: composing retrieval, a structured response format, a
symbolic verifier, and a carefully ordered training curriculum---each
individually justified---into one working system. The central question this paper
addresses is which of those components actually earn their place at 4B scale, and
at what compute cost.

Scientific reasoning presents a combined challenge. Physics problems require
formula selection, symbolic manipulation, and unit-consistent numerical
computation, whereas logic problems require quantifier handling, propositional
inference, and the detection of common fallacies. The two modalities favour
different inductive biases, yet a single model is expected to address both.

\dataset{} is the 2nd International XAI Challenge for Transparent Educational
Question-Answering, organised as an IEEE~IJCNN~2026 competition
\citep{EXACT2026}. It succeeds the XAI Challenge 2025 on explainable educational
question answering over academic regulations \citep{Nguyen2025XAIChallenge} and
broadens the scope from academic-policy logic to STEM physics, so that a system
must address both logic-based educational queries and physics problems on
electric circuits and electrostatics. The task supplies
1{,}945 training examples (1{,}213 physics, 732 logic) and a held-out evaluation
set, with a reported public baseline of 38.71\%. This study uses the challenge
only as a controlled testbed: the aim is not a competition-specific submission but
a generalisable, low-cost expert reasoning system whose design choices are meant
to transfer beyond \dataset{}. Concretely, it quantifies how far a
\basemodelname{} model can be advanced using publicly available
knowledge-engineering and fine-tuning techniques under a fixed compute budget. To
support an unambiguous interpretation, every system in this study, fine-tuned and
baseline, is evaluated under a single greedy decoding pass without any
answer-key-dependent regeneration.

The contributions of this paper are the following.
\begin{itemize}
\item A low-cost, retrieval-augmented expert reasoning system that integrates a
three-stage curriculum (general SFT $\rightarrow$ logic curriculum SFT
$\rightarrow$ GRPO reward shaping) with retrieval-augmented few-shot inference and
symbolic verification, for mixed-domain physics--logic question answering with a
4B open model, at a full end-to-end cost below US\$64 on a single A100-80GB.
\item An attribution of accuracy gains across pipeline stages under a uniform
single-pass protocol: a controlled decomposition attributes $+9.22$~pp to
retrieval-augmented exemplars ($p<10^{-3}$) and a further, borderline $+5.07$~pp
to supervised fine-tuning ($p=0.054$), while the logic-curriculum and GRPO
increments ($+1.84$~pp and $+1.38$~pp) are not individually significant at the
validation-set size; these uncertainties are reported explicitly.
\item A fair multi-prompt baseline re-evaluation in which each 7B-class model is
evaluated under its native answer format, which separates prompt-format
compliance from reasoning capability.
\item An external-benchmark evaluation on MMLU physics and FOLIO logic that
quantifies the transfer of in-domain fine-tuning, a retrieval-controlled
re-evaluation of the open baselines, and a tool-augmented evaluation that
quantifies the effect of Python and Z3 assistance.
\item A characterisation of negative results (self-consistency voting and
Dual-LoRA expert routing) and a structured error analysis that separates
answer-extraction failures from substantive reasoning errors.
\end{itemize}

The remainder of the paper is organised as follows.
Section~\ref{sec2} reviews related work. Section~\ref{sec3} formalises the task.
Section~\ref{sec4} describes the \modelname{} pipeline. Section~\ref{sec5}
presents the experimental setup. Section~\ref{sec6} reports the results.
Section~\ref{sec7} discusses limitations, and Section~\ref{sec8} concludes.

%% ============================================================
\section{Related Work}\label{sec2}
%% ============================================================

\subsection{Chain-of-thought and self-consistency}
\citet{Wei2022} showed that step-by-step prompting improves LLM performance on
arithmetic and symbolic reasoning, and \citet{Kojima2022} demonstrated
zero-shot chain-of-thought prompting. \citet{Wang2023sc} proposed
self-consistency decoding, which samples multiple reasoning paths and aggregates
the most frequent answer. The present results indicate that self-consistency
provides no net benefit at 4B scale on this task (Section~\ref{sec6}).

\subsection{Mathematics and science language models}
DeepSeekMath \citep{Shao2024} combined large-scale mathematical pre-training
with GRPO. Qwen2.5-Math \citep{Yang2024} and the broader Qwen family
\citep{Qwen2025} extended this approach across model scales. Llemma
\citep{Azerbayev2024} performed continued pre-training on scientific and proof
corpora. MetaMath \citep{Yu2023} and MAmmoTH \citep{Yue2023} constructed
augmented mathematical reasoning datasets. The present work operates at a
smaller scale (4B) in a two-domain setting (physics and logic) and introduces a
unit-consistency reward term.

\subsection{Neuro-symbolic integration}
Logic-LM \citep{Pan2023} and LINC \citep{Olausson2023} couple LLMs with formal
solvers for logical reasoning, and PAL \citep{Gao2022} delegates arithmetic to a
Python interpreter. Math-Shepherd \citep{Wang2024mathshepherd} and
process-supervision methods \citep{Lightman2023} provide step-level reward
labels. This study integrates the Z3 SMT solver into the correctness reward and
into a tool-augmented inference variant, while relying on outcome-level rewards
to avoid step-level annotation.

\subsection{Parameter-efficient fine-tuning and reinforcement learning}
LoRA \citep{Hu2022} decomposes weight updates into low-rank matrices, and QLoRA
\citep{Dettmers2023} extends it to quantised backbones; LoRAHub
\citep{Huang2023} composes multiple domain adapters. InstructGPT
\citep{Ouyang2022} established RLHF via PPO \citep{Schulman2017}; Direct
Preference Optimization \citep{Rafailov2023} removes the explicit reward model;
and GRPO \citep{Shao2024} computes relative rewards within a generation group,
which makes single-GPU training of 4B-class models practical. Curriculum
learning \citep{Bengio2009} motivates the staged data ordering used here.

\subsection{Explainable educational QA challenges}
The closest predecessor to \dataset{} is the XAI Challenge 2025, an explainable
educational question-answering competition organised by the URA Research Group at
Ho Chi Minh City University of Technology together with the TRNS-AI workshop at
IJCNN~2025 \citep{Nguyen2025XAIChallenge}. That challenge required systems to
produce not only final answers but also logic-grounded natural-language
explanations for questions about university policies; its dataset (481 training
and 50 test records) was constructed from first-order-logic templates with Z3
validation and expert-student review, with premises supplied in both natural
language and first-order logic. \dataset{} continues this line of work and
broadens it from academic-regulation logic to STEM physics, evaluating
submissions along answer correctness, explanation quality, and reasoning depth
while encouraging---but not requiring---symbolic tools such as Z3
\citep{EXACT2026}. In contrast to the 2025 challenge report, which presents the
competition, dataset, and evaluation protocol, the present paper is a
system-level empirical study: it fixes a single 4B open model and reports
controlled ablations of retrieval, supervised fine-tuning, logic-curriculum
adaptation, GRPO reward shaping, tool augmentation, and self-consistency.
We do not propose a new training algorithm; the contribution is a controlled
attribution of which standard components actually pay off at 4B scale and which
do not, isolated by the retrieval-only control (cfg0-R) and a uniform single-pass
protocol.
Table~\ref{tabrelated} situates \modelname{} relative to representative reasoning
systems and to these two challenge artefacts.

\begin{table*}[t]
\centering
\caption{Positioning of \modelname{} relative to representative LLM reasoning
systems and to the EXACT challenge artefacts. ``Tuning'' indicates whether the
approach trains model weights.\label{tabrelated}}
\begin{tabular*}{\textwidth}{@{\extracolsep\fill}p{0.21\textwidth}p{0.18\textwidth}p{0.06\textwidth}p{0.40\textwidth}@{}}
\toprule
\textbf{Work} & \textbf{Focus / domain} & \textbf{Tuning} & \textbf{Relation to \modelname{}} \\
\midrule
XAI Challenge 2025 report \citep{Nguyen2025XAIChallenge} & Academic-regulation QA & --- & Predecessor; defines the challenge and dataset, not a system study \\
\dataset{} task \citep{EXACT2026} & Logic $+$ STEM physics & --- & Official task, rules, and evaluation; not a method \\
CoT / self-consistency \citep{Wei2022,Wang2023sc} & Prompt-time reasoning & No & Evaluated after 4B adaptation; no net gain observed \\
PAL \citep{Gao2022} & Program-aided arithmetic & No & Python assistance tested; only marginal gain \\
Logic-LM / LINC \citep{Pan2023,Olausson2023} & Logic with solvers & No & Z3 used, but the NL-logic parser limits coverage \\
DeepSeekMath / Qwen-Math \citep{Shao2024,Yang2024} & Maths-specialised LLMs & Yes & Motivation and baselines for low-cost GRPO \\
LoRA / QLoRA \citep{Hu2022,Dettmers2023} & Parameter-efficient tuning & Yes & Applied to a 4B model under a fixed budget \\
\modelname{} (this work) & Logic $+$ physics & Yes & Retrieval $+$ SFT $+$ GRPO with reproducible ablations \\
\bottomrule
\end{tabular*}
\end{table*}

%% ============================================================
\section{Problem Formulation}\label{sec3}
%% ============================================================

Let $\mathcal{D}=\mathcal{D}_{\text{phys}}\cup\mathcal{D}_{\text{logic}}$ denote
the \dataset{} dataset \citep{EXACT2026}. Each sample
$(q_i,a_i^{*},s_i)\in\mathcal{D}$ comprises a
question $q_i$, a ground-truth answer $a_i^{*}$, and a domain label
$s_i\in\{\text{physics},\text{logic}\}$. A model $M$ with parameters $\theta$
maps a question to a free-text response $\hat{y}_i=M_\theta(q_i)$. A structured
extraction operator recovers the candidate answer $a_i=\texttt{extract}
(\hat{y}_i)$, and accuracy is
\begin{equation}
\mathrm{Acc}=\frac{1}{|\mathcal{D}_{\text{val}}|}
\sum_{i=1}^{|\mathcal{D}_{\text{val}}|}
\mathbf{1}\!\left[\mathrm{verify}(a_i,a_i^{*},s_i,q_i)\right],
\end{equation}
where $\mathrm{verify}(\cdot)$ applies Z3-based symbolic checking, when a
parseable formal representation is available, and SI-unit-normalised numerical
comparison for physics samples. Each response is produced by a single greedy
decoding pass; no component of the evaluation uses the ground-truth label to
decide whether to regenerate a response. The official challenge scores
submissions along three dimensions---answer correctness, explanation quality, and
reasoning depth---and encourages optional structured evidence such as
first-order-logic derivations and chain-of-thought steps \citep{EXACT2026}; the
present study targets the answer-correctness dimension and additionally examines
answer formatting and reasoning errors in Section~\ref{sec6}.

%% ============================================================
\section{Methodology}\label{sec4}
%% ============================================================

Figure~\ref{fig1} summarises the pipeline. All stages apply LoRA
($r{=}32$, $\alpha{=}64$) to \basemodelname{}.

\begin{figure}[t]
\centerline{\includegraphics[width=\linewidth]{fig1_pipeline}}
\caption{Overview of the \modelname{} training and inference pipeline. Stage~1
performs general supervised fine-tuning; Stage~2 applies a logic-focused
curriculum stage; Stage~3 applies GRPO with a composite reward. At inference,
the fine-tuned configurations retrieve three in-context exemplars. Optional
Dual-LoRA routing and self-consistency voting are evaluated as additional
configurations.\label{fig1}}
\end{figure}

\subsection{Structured response format}
All models are instructed to produce a three-tag response,
$\langle\texttt{reasoning}\rangle\ldots\langle/\texttt{reasoning}\rangle\,
\langle\texttt{answer}\rangle\ldots\langle/\texttt{answer}\rangle\,
\langle\texttt{explanation}\rangle\ldots\langle/\texttt{explanation}\rangle$,
which supports reliable answer extraction and enables a format-based reward
term.

\subsection{Phase 1: supervised fine-tuning}
\basemodelname{} is fine-tuned on the full training set (1{,}945 samples) with a
cross-entropy loss on the target response. The configuration is three epochs,
per-device batch size four with gradient accumulation four (effective batch
16), learning rate $2\times10^{-4}$, LoRA $r{=}32$, $\alpha{=}64$, dropout
$0.05$, applied to all seven projection matrices. Training completes in
58.8~min and converges to a training loss of $0.3085$.

\subsection{Phase 1.5: logic curriculum fine-tuning}
A second SFT stage is restricted to the logic subset (732 samples) with a
reduced learning rate of $5\times10^{-5}$ for two epochs. This curriculum stage
\citep{Bengio2009} allocates additional capacity to formal reasoning, which is
under-represented (37.6\% of the training data). Training completes in 15.7~min
and converges to a loss of $0.1326$.

\subsection{Phase 2: GRPO with reward shaping}
For a completion $y$ with ground truth $a^{*}$ and domain $s$, the reward is
\begin{equation}
\begin{split}
R(y,a^{*},s)={}&w_f\,r_{\text{format}}(y)+w_c\,r_{\text{correct}}(y,a^{*},s)\\
&{}+w_u\,r_{\text{unit}}(y,a^{*})+\lambda\,\mathbf{1}[\,|y_{\text{reason}}|>800\,],
\end{split}
\label{eq:reward}
\end{equation}
with $w_f{=}0.30$, $w_c{=}0.60$, $w_u{=}0.10$, and $\lambda{=}-0.10$, where
$|y_{\text{reason}}|$ is the reasoning token count. The format term awards
$0.10$ per present tag (maximum $0.30$). The correctness term is binary and
applies a $\le 5\%$ relative tolerance after SI-unit normalisation for physics
and normalised string or Z3-verified matching for logic. The unit-consistency
term equals one when the extracted answer carries a unit dimensionally
compatible with the ground-truth unit. GRPO \citep{Shao2024} is applied from the
Phase-1.5 checkpoint with $K{=}4$ completions per prompt; the group-normalised
advantage is $\hat{r}_k=(R_k-\bar{R})/(\sigma_R+\epsilon)$, with KL coefficient
$\beta{=}0.04$, learning rate $10^{-6}$, and 250 steps. Training completes in
79.8~min and reaches a loss of $0.0295$.

\subsection{Retrieval-augmented inference and optional components}
At inference, the fine-tuned configurations prepend three in-context exemplars
retrieved from the training set by dense similarity to the query; the zero-shot
configuration uses no exemplars. As additional configurations, two LoRA
adapters (one physics-specialised, one logic-specialised) are routed by a
TF-IDF and logistic-regression classifier (5{,}000 n-gram features); when the
router confidence exceeds $0.80$ the corresponding adapter is selected,
otherwise a keyword heuristic is used. Self-consistency voting with $N{=}5$
samples is applied as a post-processing step in the full-agent configuration.

%% ============================================================
\section{Experimental Setup}\label{sec5}
%% ============================================================

\subsection{Dataset}
Table~\ref{tab1} summarises the data split. The raw sources are a logic corpus
and a physics problem set, combined and divided by a 90/10 stratified
train/validation split.

\begin{table}[t]
\centering
\caption{Dataset statistics for \dataset{}.\label{tab1}}
\begin{tabular*}{\columnwidth}{@{\extracolsep\fill}lrrr@{\extracolsep\fill}}
\toprule
\textbf{Split} & \textbf{Total} & \textbf{Physics} & \textbf{Logic} \\
\midrule
Train & 1{,}945 & 1{,}213 & 732 \\
Validation & 217 & 141 & 76 \\
\midrule
Total & 2{,}162 & 1{,}354 & 808 \\
\bottomrule
\end{tabular*}
\end{table}

\textit{Train--validation overlap.} Because both splits derive from a single
released corpus, we audited the validation set for leakage with three
complementary measures: the TF-IDF cosine similarity of each validation question
against the full training set, a character $8$-gram Jaccard overlap that (unlike
word-level features) preserves the numerical values distinguishing otherwise
identically worded problems, and inspection of the top-$3$ retrieval exemplars
that cfg0-R actually uses. Surface similarity is high---$66$ of $217$ validation
questions reach a maximum TF-IDF cosine $\geq 0.90$---but this overwhelmingly
reflects template reuse: the corpus instantiates the same physics or logic
schema with different numerical values, so near-identical wording maps to a
different answer. Under character-level near-identity ($8$-gram
Jaccard $\geq 0.95$), only $6$ of $217$ questions are genuine near-duplicates,
and $5$ ($2.3\%$) are both near-identical and share the gold answer; three
questions appear verbatim in training, one of them carrying a different gold
label (a numerical-versus-categorical mismatch of the kind documented in the
error analysis). This residual ${\approx}2\%$ overlap is far too small to account
for the $17.51$~pp headline gain, and the retrieved exemplars are verbatim copies
of the validation question in only these three cases (mean top-exemplar cosine
$0.68$). The audit script (\texttt{step10\_leakage\_check.py}) is released with
the code.

\subsection{Baselines and evaluation}
Five open-weight 7B-class models (Qwen2.5-Math-7B-Instruct \citep{Yang2024},
Qwen2-Math-7B-Instruct, Llemma-7B \citep{Azerbayev2024},
Mistral-7B-Instruct-v0.3, and DeepSeek-Math-7B-Instruct) and one proprietary
model (Gemini-2.5-flash-lite, accessed through an API) are evaluated on the 217
validation samples. The challenge rules prohibit closed-source models in
submitted systems \citep{EXACT2026}; the Gemini baseline is therefore reported
only as a non-submission diagnostic reference and is not part of the
\modelname{} system itself, all components of which use
open-weight models for reproducibility. To separate format compliance from reasoning capability,
each open baseline is evaluated under a fair multi-prompt protocol that tries
three answer formats (XML tags, chain-of-thought with a boxed answer, and
chain-of-thought with a plain answer line) and reports the best per model. The
primary metric is overall accuracy, with physics and logic sub-accuracies
reported separately. Statistical significance between configuration pairs is
assessed by McNemar's test (two-sided, continuity-corrected), and Wilson score
intervals are reported for overall accuracy. Decoding uses greedy generation
with a maximum of 1{,}024 new tokens for all systems. Table~\ref{tab2} lists the
principal hyper-parameters.

\begin{table}[t]
\centering
\caption{Principal hyper-parameters.\label{tab2}}
\begin{tabular*}{\columnwidth}{@{\extracolsep\fill}ll@{\extracolsep\fill}}
\toprule
\textbf{Parameter} & \textbf{Value} \\
\midrule
Base model & \basemodelname{} (BF16) \\
LoRA $r$ / $\alpha$ / dropout & 32 / 64 / 0.05 \\
SFT learning rate / epochs & $2\times10^{-4}$ / 3 \\
Logic SFT learning rate / epochs & $5\times10^{-5}$ / 2 \\
GRPO learning rate / steps & $10^{-6}$ / 250 \\
GRPO KL coefficient $\beta$ & 0.04 \\
GRPO completions per prompt $K$ & 4 \\
Reward weights $(w_f,w_c,w_u)$ & $(0.30,0.60,0.10)$ \\
Length penalty $\lambda$ & $-0.10$ if ${>}800$ tokens \\
Retrieval exemplars (fine-tuned) & 3 \\
Max new tokens & 1{,}024 \\
GPU & A100-SXM4-80GB \\
\bottomrule
\end{tabular*}
\end{table}

%% ============================================================
\section{Results and Analysis}\label{sec6}
%% ============================================================

\subsection{Main results}
Table~\ref{tab3} reports the comparison, and Figure~\ref{fig2} visualises the
per-stage accuracy. The strongest configuration, cfg3 (SFT, logic SFT, and
GRPO), reaches 47.47\% overall accuracy (its increment over the curriculum stage
being within validation-set noise), 17.51~pp above the \basemodelname{}
zero-shot configuration and above all six baselines (Figure~\ref{fig3}). The
fine-tuned configurations use retrieval-augmented exemplars, whereas the
zero-shot configuration and the baselines do not. To separate the two effects,
configuration cfg0-R applies the same three-exemplar retrieval to the base model
without any fine-tuning: retrieval alone accounts for $+9.22$~pp
(cfg0$\rightarrow$cfg0-R, $p<10^{-3}$), and fine-tuning adds a further $+5.07$~pp
(cfg0-R$\rightarrow$cfg1, $p=0.054$; Table~\ref{tab4}). The best configuration
also exceeds the reported \dataset{}
public baseline of 38.71\% by 8.76~pp; because that baseline is not reproduced
under the present single-pass protocol, it is used only as a contextual
reference, and the controlled comparisons rely on cfg0 and the fair 7B-class
baselines. Although cfg3 attains the highest overall accuracy, its improvement
over cfg2 is driven by physics: logic accuracy ($18.42\%$) does not improve over
cfg2 ($21.05\%$), which indicates that the GRPO reward is more effective for
answer formatting and numerical reasoning than for formal logical inference. This
$2.63$~pp difference corresponds to two of 76 logic items and lies within the
non-significant cfg2$\rightarrow$cfg3 transition ($p=0.61$, Table~\ref{tab4}); it
should therefore be read as GRPO leaving logic accuracy unchanged rather than as
a genuine regression. The
low absolute scores of the maths-specialised 7B baselines should be read in light
of the mixed physics--logic format and the strict answer-extraction protocol
(Appendix~\ref{appA}) rather than as a general ranking of those models: even
under the fair best-of-three protocol, the heterogeneous answer space of
\dataset{} penalises systems that are not adapted to it.

\begin{table*}[t]
\centering
\caption{Main results on the \dataset{} validation set (217 samples: 141
physics, 76 logic) under a single greedy decoding pass. Baselines report the
best of three answer formats. The double dagger ($\ddagger$) marks McNemar
$p<0.01$ versus cfg0 for the two transitions tested in Table~\ref{tab4}
(cfg0$\rightarrow$cfg1, cfg0$\rightarrow$cfg3).\label{tab3}}
\begin{tabular*}{\textwidth}{@{\extracolsep\fill}llrrr@{\extracolsep\fill}}
\toprule
& \textbf{System} & \textbf{Overall} & \textbf{Physics} & \textbf{Logic} \\
\midrule
\multicolumn{5}{@{}l}{\textit{Zero-shot baselines (best-of-three prompt)}} \\
B1 & Qwen2-Math-7B-Instruct & 29.03 & 35.46 & 17.11 \\
B2 & Qwen2.5-Math-7B-Instruct & 24.88 & 34.75 & \hphantom{0}6.58 \\
B3 & Gemini-2.5-flash-lite & 24.88 & 32.62 & 10.53 \\
B4 & Mistral-7B-Instruct-v0.3 & 16.59 & 13.48 & 22.37 \\
B5 & DeepSeek-Math-7B-Instruct & 12.44 & 14.89 & \hphantom{0}7.89 \\
B6 & Llemma-7B & 11.52 & \hphantom{0}4.96 & 23.68 \\
\midrule
\multicolumn{5}{@{}l}{\textit{Ours (fine-tuned \basemodelname{})}} \\
cfg0 & Zero-shot & 29.95 & 38.30 & 14.47 \\
cfg0-R & $+$ Retrieval (no fine-tuning)$^{\ddagger}$ & 39.17 & 50.35 & 18.42 \\
cfg1 & $+$ SFT$^{\ddagger}$ & 44.24 & 58.16 & 18.42 \\
cfg2 & $+$ Logic SFT & 46.08 & 59.57 & 21.05 \\
cfg3 & $+$ GRPO$^{\ddagger}$ & \textbf{47.47} & \textbf{63.12} & 18.42 \\
cfg4 & Dual-LoRA $+$ Router & 43.78 & 55.32 & 22.37 \\
cfg5 & Full agent (SC$\times5$) & 44.70 & 56.74 & 22.37 \\
\bottomrule
\end{tabular*}
\end{table*}

\begin{figure}[t]
\centerline{\includegraphics[width=\linewidth]{fig_stage_accuracy}}
\caption{Overall, physics, and logic accuracy across the pipeline stages under a
single greedy decoding pass; cfg0-R (shaded) is the retrieval-only control (base
model with retrieval, no fine-tuning). The dashed line marks the best overall
accuracy (cfg3, 47.47\%). Physics accuracy increases monotonically through cfg3,
whereas logic accuracy remains below 23\% for all configurations.\label{fig2}}
\end{figure}

\begin{figure}[t]
\centerline{\includegraphics[width=\linewidth]{fig_vs_baselines}}
\caption{Overall accuracy of the fine-tuned 4B model (cfg3) against the six
zero-shot baselines under the fair best-of-three protocol. The fine-tuned model
exceeds the strongest baseline by 18.4~pp.\label{fig3}}
\end{figure}

\subsection{Does retrieval explain the baseline gap?}
A natural concern is whether the margin over the 7B baselines merely reflects the
retrieval exemplars available to the fine-tuned configurations. To test this, the
same three-exemplar retrieval used by cfg0-R is supplied to each of the five
open-weight baselines under the identical best-of-three protocol
(Table~\ref{tabretr}). Retrieval helps the baselines only marginally on average
($+2.6$~pp): the strongest baseline rises from $29.03\%$ to $30.88\%$, and only
DeepSeek-Math gains substantially ($+9.22$~pp, from a low base), while one
baseline is unchanged and another declines slightly. Crucially, even with
identical retrieval the best baseline ($30.88\%$) remains $8.29$~pp below the
retrieval-only control cfg0-R ($39.17\%$) and $16.59$~pp below cfg3 ($47.47\%$).
The advantage of \modelname{} is therefore not attributable to retrieval access:
the curriculum-fine-tuned 4B model exploits the same exemplars far more
effectively than larger maths-specialised 7B models, consistent with the
format-alignment and answer-extraction analysis above.

\begin{table}[t]
\centering
\caption{Effect of supplying the cfg0-R retrieval exemplars to the five
open-weight baselines (best of three prompt formats, $n=217$). Even with
identical retrieval, the strongest baseline trails the retrieval-only control
cfg0-R and the full pipeline cfg3.\label{tabretr}}
\begin{tabular*}{\columnwidth}{@{\extracolsep\fill}lrrr@{\extracolsep\fill}}
\toprule
\textbf{Baseline} & \textbf{No retr.} & \textbf{$+$Retr.} & \textbf{$\Delta$} \\
\midrule
Qwen2-Math-7B    & 29.03 & 30.88 & $+1.85$ \\
Qwen2.5-Math-7B  & 24.88 & 27.65 & $+2.77$ \\
Mistral-7B       & 16.59 & 16.59 & $+0.00$ \\
DeepSeek-Math-7B & 12.44 & 21.66 & $+9.22$ \\
Llemma-7B        & 11.52 & 10.60 & $-0.92$ \\
\midrule
\textit{cfg0-R} (ours, retrieval only) & --- & \textit{39.17} & \\
\textit{cfg3} (ours, full pipeline)    & --- & \textit{47.47} & \\
\bottomrule
\end{tabular*}
\end{table}

\subsection{Statistical confidence and significance}
Table~\ref{tab4} reports Wilson score intervals and pairwise McNemar tests
computed on the single-pass per-sample predictions. The zero-shot-to-SFT
(cfg0$\rightarrow$cfg1) and zero-shot-to-GRPO (cfg0$\rightarrow$cfg3)
transitions are both significant; the logic-SFT increment
(cfg1$\rightarrow$cfg2, $p=0.39$), the GRPO-over-SFT increment
(cfg2$\rightarrow$cfg3, $p=0.61$), and the self-consistency change
(cfg3$\rightarrow$cfg5, $p=0.24$) are not significant at the $0.05$ level. The
overlap among the cfg1--cfg5 intervals indicates that differences among the
fine-tuned configurations are small relative to the validation-set size. Because
cfg1 includes both supervised fine-tuning and retrieval-augmented exemplars, the
cfg0-to-cfg1 comparison should not be read as the isolated effect of
fine-tuning; the cfg0-R control separates retrieval from fine-tuning and shows
that retrieval accounts for the larger share of the first-stage improvement.

\begin{table*}[t]
\centering
\caption{Wilson 95\% intervals (left) and pairwise McNemar tests (right) under
the single-pass protocol ($n=217$). For a pair $A\rightarrow B$, $b$ is the
count where $B$ is correct and $A$ is wrong, and $c$ is the reverse.\label{tab4}}
\begin{tabular*}{\textwidth}{@{\extracolsep\fill}lr@{\hskip 24pt}lrrrr@{\extracolsep\fill}}
\toprule
\textbf{System} & \textbf{Overall (95\% CI)} & \textbf{Pair} & \textbf{$\Delta$Acc} & \textbf{$b$} & \textbf{$c$} & \textbf{$p$} \\
\midrule
cfg0 & 29.95 [24.3, 36.4] & cfg0$\rightarrow$cfg0-R & $+9.22$ & 26 & \hphantom{0}6 & $7.8\times10^{-4}$ \\
cfg0-R & 39.17 [32.9, 45.8] & cfg0-R$\rightarrow$cfg1 & $+5.07$ & 19 & \hphantom{0}8 & $0.054$ \\
cfg1 & 44.24 [37.8, 50.9] & cfg0$\rightarrow$cfg3 & $+17.51$ & 43 & \hphantom{0}5 & $9.3\times10^{-8}$ \\
cfg2 & 46.08 [39.6, 52.7] & cfg1$\rightarrow$cfg2 & $+1.84$ & \hphantom{0}8 & \hphantom{0}4 & $0.39$ \\
cfg3 & 47.47 [40.9, 54.1] & cfg2$\rightarrow$cfg3 & $+1.38$ & \hphantom{0}9 & \hphantom{0}6 & $0.61$ \\
cfg4 & 43.78 [37.3, 50.4] & cfg3$\rightarrow$cfg5 & $-2.76$ & \hphantom{0}6 & 12 & $0.24$ \\
cfg5 & 44.70 [38.2, 51.4] & & & & & \\
\bottomrule
\end{tabular*}
\end{table*}

\subsection{Multi-seed reliability}
Table~\ref{tab5} reports cfg3 re-evaluated with three generation seeds. Under the
single-pass protocol the overall standard deviation is 0.70~pp, which is small
relative to the 17.51~pp zero-shot-to-cfg3 gain and confirms that the reported
accuracy is stable across seeds.

\begin{table}[t]
\centering
\caption{cfg3 multi-seed overall accuracy (single greedy pass).\label{tab5}}
\begin{tabular*}{\columnwidth}{@{\extracolsep\fill}lr@{\extracolsep\fill}}
\toprule
\textbf{Seed} & \textbf{Overall accuracy (\%)} \\
\midrule
42 & 47.47 \\
1337 & 46.08 \\
2024 & 46.54 \\
\midrule
Mean $\pm$ std & $46.70 \pm 0.70$ \\
\bottomrule
\end{tabular*}
\end{table}

\subsection{External generalisation: physics and logic}
To assess transfer beyond the in-domain split, cfg0, cfg2, and cfg3 are evaluated
on the MMLU \citep{Hendrycks2021} high-school and college physics subsets (253
four-option multiple-choice questions) without additional fine-tuning
(Table~\ref{tab6}). The fine-tuned configurations improve over the zero-shot
model by up to 6.33~pp, with monotonic improvement across stages, which provides
an external sanity check that the in-domain curriculum transfers rather than
overfitting to the \dataset{} format. These MMLU results serve as a transfer
check rather than a competitive benchmark claim.

\begin{table}[t]
\centering
\caption{Zero-shot transfer to MMLU physics (253 questions: 151 high-school, 102
college).\label{tab6}}
\begin{tabular*}{\columnwidth}{@{\extracolsep\fill}lrrr@{\extracolsep\fill}}
\toprule
\textbf{Configuration} & \textbf{Overall} & \textbf{HS} & \textbf{College} \\
\midrule
cfg0 (zero-shot) & 42.29 & 41.06 & 44.12 \\
cfg2 ($+$ logic SFT) & 47.04 & 47.68 & 46.08 \\
cfg3 ($+$ GRPO) & \textbf{48.62} & 46.36 & 51.96 \\
\bottomrule
\end{tabular*}
\end{table}

The same transfer test is repeated for logic on FOLIO \citep{Han2022folio}, a
first-order-logic entailment benchmark whose three-way label space
(entailed / contradicted / undetermined) matches the \dataset{} logic format
(Table~\ref{tabfolio}). The zero-shot base model scores at the majority-class
level ($34.48\%$), whereas the logic-curriculum configuration cfg2 reaches
$48.28\%$---a $+13.80$~pp gain---and cfg3 attains $46.80\%$. The logic curriculum
therefore induces genuine, transferable logical reasoning rather than overfitting
to the \dataset{} answer format, and the small cfg2$\rightarrow$cfg3 decrease
($-1.48$~pp) mirrors the in-domain result that GRPO does not improve logic. The
larger logic transfer ($+13.80$~pp) than physics transfer ($+6.33$~pp) indicates
that the curriculum's logic stage, not only its physics exposure, carries beyond
the \dataset{} split; this is, however, a larger \emph{relative} gain from a
near-chance base, and in-domain logic accuracy remains below 23\%, so formal
logic stays the principal bottleneck.

\begin{table}[t]
\centering
\caption{Zero-shot transfer to FOLIO logic (203 validation items;
majority-class baseline $\approx 35.5\%$).\label{tabfolio}}
\begin{tabular*}{\columnwidth}{@{\extracolsep\fill}lrr@{\extracolsep\fill}}
\toprule
\textbf{Configuration} & \textbf{Accuracy} & \textbf{$\Delta$ vs cfg0} \\
\midrule
cfg0 (zero-shot)     & 34.48 & --- \\
cfg2 ($+$ logic SFT) & \textbf{48.28} & $+13.80$ \\
cfg3 ($+$ GRPO)      & 46.80 & $+12.32$ \\
\bottomrule
\end{tabular*}
\end{table}

\subsection{Tool-augmented reasoning}
Table~\ref{tab7} compares direct inference with a tool-augmented variant of the
best checkpoint, in which a Python evaluator handles physics arithmetic and the
Z3 solver handles logic inference; the model answer is retained when the solver
abstains. Both arms use a single-pass protocol \emph{without} retrieval
exemplars, so the absolute values are not directly comparable to
Table~\ref{tab3}; the purpose is to isolate the tool effect. Tool augmentation
yields a marginal overall change ($+0.46$~pp), with a small physics gain from the
Python evaluator and a small logic decrease. The Z3 solver abstained on the
natural-language logic items because their premises were not reliably parseable
into a formal specification, which limits the benefit of symbolic tooling under
the current parser.

\begin{table}[t]
\centering
\caption{Direct versus tool-augmented inference on the best checkpoint
($n=217$, single pass, no retrieval exemplars). Tool activations: Python
evaluator on 62 samples; Z3 abstained on the remaining cases.\label{tab7}}
\begin{tabular*}{\columnwidth}{@{\extracolsep\fill}lrrr@{\extracolsep\fill}}
\toprule
\textbf{Arm} & \textbf{Overall} & \textbf{Physics} & \textbf{Logic} \\
\midrule
A: direct & 39.63 & 49.65 & 21.05 \\
B: $+$ Python / Z3 tools & 40.09 & 51.06 & 19.74 \\
\midrule
$\Delta$ (B $-$ A) & $+0.46$ & $+1.41$ & $-1.31$ \\
\bottomrule
\end{tabular*}
\end{table}

\subsection{Reward component ablation}
Table~\ref{tab8} reports four reward variants trained for 150 GRPO steps from the
Phase-1 checkpoint and evaluated under single-pass decoding without retrieval.
The correctness-only variant (R4) attains the highest accuracy at this budget,
1.84~pp above the full mixture (R1); removing the format term (R3) reduces
physics accuracy by 2.84~pp. At this 150-step budget the auxiliary
unit-consistency and format terms do not provide a net accuracy benefit relative
to a correctness-only objective. A plausible explanation is that the correctness
reward supplies the most direct training signal and that auxiliary terms add
gradient variance that is not amortised within 150 steps; whether the ordering
reverses at longer horizons was not evaluated within the compute budget.

\begin{table}[t]
\centering
\caption{Reward component ablation (150 GRPO steps from the Phase-1 checkpoint,
single pass, no retrieval, $n=217$).\label{tab8}}
\begin{tabular*}{\columnwidth}{@{\extracolsep\fill}llrr@{\extracolsep\fill}}
\toprule
& \textbf{Reward configuration} & \textbf{Overall} & \textbf{Physics} \\
\midrule
R1 & Full (format, correct, unit, length) & 39.63 & 49.65 \\
R2 & No unit term & 39.17 & 49.65 \\
R3 & No format term & 38.25 & 46.81 \\
R4 & Correctness only & \textbf{41.47} & \textbf{51.06} \\
\bottomrule
\end{tabular*}
\end{table}

\subsection{Compute cost}
Table~\ref{tab9} summarises the compute profile. Phase-1 SFT provides the highest
accuracy per training dollar, and the full pipeline cost is below one A100-day.

\begin{table}[t]
\centering
\caption{Compute profile on one A100-SXM4-80GB at US\$3.67 per GPU-hour.\label{tab9}}
\begin{tabular*}{\columnwidth}{@{\extracolsep\fill}lrr@{\extracolsep\fill}}
\toprule
\textbf{Stage} & \textbf{Time (min)} & \textbf{Cost (US\$)} \\
\midrule
SFT Phase 1 & 58.8 & 3.60 \\
Logic SFT (Phase 1.5) & 15.7 & 0.96 \\
GRPO Phase 2 (mixed) & 79.8 & 4.89 \\
GRPO Phase 2 (physics) & 76.6 & 4.69 \\
GRPO Phase 2 (logic) & 62.9 & 3.85 \\
\midrule
Total training & 293.8 & \textbf{18.0} \\
\bottomrule
\end{tabular*}
\end{table}

\subsection{Error analysis}
Because the canonical evaluation retained only resolved outputs, a dedicated
single-pass run of cfg3 was performed to retain the full generations required
for error classification. Table~\ref{tab10} reports a rule-based error taxonomy
on that run ($n=117$ incorrect predictions; the run scores 46.08\% overall,
within the seed variation of Table~\ref{tab5}). The largest category is
answer-extraction failure (72.6\%): the response does not yield an answer that
the automated verifier can reconcile with the ground truth, and inspection
indicates that many of these are incomplete generations that terminate before a
final answer is emitted. The second category is logical-fallacy and
quantifier error (17.9\%), concentrated in the logic subset. Two observations
follow. First, because a large fraction of measured errors are extraction or
generation-length failures rather than substantive reasoning errors, part of the
measured error may be attributable to extraction limitations rather than to
reasoning; a manual annotation protocol is required to quantify this fraction,
which is left to future work. Second, among substantive reasoning errors the
logical-fallacy category dominates, consistent with the persistent
physics--logic accuracy gap and identifying formal logic as the principal
bottleneck at this scale.

\begin{table}[t]
\centering
\caption{Rule-based error taxonomy for incorrect cfg3 predictions ($n=117$,
single-pass run). Labels are produced by an automated classifier; manual
validation is left to future work.\label{tab10}}
\begin{tabular*}{\columnwidth}{@{\extracolsep\fill}llrr@{\extracolsep\fill}}
\toprule
\textbf{ID} & \textbf{Error type} & \textbf{Count} & \textbf{\%} \\
\midrule
E6 & Answer-extraction / incomplete output & 85 & 72.6 \\
E4 & Logical fallacy / quantifier error & 21 & 17.9 \\
E0 & Unclassified & \hphantom{0}9 & \hphantom{0}7.7 \\
E1 & Unit or dimension error & \hphantom{0}2 & \hphantom{0}1.7 \\
\bottomrule
\end{tabular*}
\end{table}

To make these categories concrete, Table~\ref{tab11} lists representative cfg3
failures drawn from the single-pass run. The dominant pattern is an
answer-extraction mismatch in which the model reasons soundly but emits an answer
whose form the verifier cannot reconcile with a categorical ground truth,
together with truncated generations that stop before a final answer is emitted;
the substantive logical failures are typically faulty quantifier or existential
reasoning.

\begin{table*}[t]
\centering
\caption{Representative cfg3 errors (single-pass run) illustrating the dominant
categories of Table~\ref{tab10}.\label{tab11}}
\begin{tabular*}{\textwidth}{@{\extracolsep\fill}p{0.16\textwidth}p{0.44\textwidth}p{0.30\textwidth}@{}}
\toprule
\textbf{Category} & \textbf{Symptom (abridged)} & \textbf{Interpretation} \\
\midrule
E6: extraction mismatch & Capacitor charge after disconnection: the model answers ``$100~\mu$C'' (the correct value), but the ground truth is the categorical label ``unchanged'' & Sound reasoning scored wrong; numerical versus categorical answer space \\
E6: incomplete output & Electric-field problem: the generation terminates mid-derivation (``\ldots\ A is at $x=1$~cm. B is'') with no final answer & Generation-length limit; constrained decoding could recover the answer \\
E4: logical fallacy & Quantifier item: the model declares the query a tautology and answers ``Yes'' where the entailment is ``No'' & Faulty existential/quantifier reasoning, not an extraction issue \\
\bottomrule
\end{tabular*}
\end{table*}

%% ============================================================
\section{Discussion}\label{sec7}
%% ============================================================

\subsection{Self-consistency at 4B scale}
Self-consistency is effective when errors are independent across samples
\citep{Wang2023sc}. The present error profile indicates that a large share of
failures are systematic, either fixed answer-format failures or repeated
quantifier errors, which majority voting reinforces rather than corrects. This
accounts for the non-significant decrease from cfg3 to cfg5
($-2.76$~pp, $p=0.24$) despite five times the inference compute.

\subsection{Domain routing}
The TF-IDF and logistic-regression router classifies clean physics and logic
items reliably, but border-zone questions, such as physics problems phrased as
logical conditionals, are routed to the wrong adapter. The Dual-LoRA
configuration (cfg4, 43.78\%) does not exceed the single mixed adapter
(cfg3, 47.47\%), which indicates that routing noise offsets the benefit of
domain specialisation at this scale.

\subsection{Limitations}
\begin{itemize}
\item \textit{Validation-set size.} With 217 samples, Wilson intervals span
roughly $\pm7$~pp; differences below approximately 5~pp are not statistically
separable, including the GRPO-over-SFT increment ($+1.38$~pp, $p=0.61$). The 217
samples are a 90/10 stratified split of the released training data; the official
competition test labels were not available, so no independent test-set result is
reported. Consequently only the large cfg0$\rightarrow$cfg3 gain
($p=9.3\times10^{-8}$) survives the small-sample caveat, whereas the increments
among the fine-tuned configurations do not.
\item \textit{Inference asymmetry.} The fine-tuned configurations use three
retrieved in-context exemplars, whereas the baselines do not. Configuration
cfg0-R (Table~\ref{tab3}) isolates the retrieval contribution from fine-tuning;
cross-family comparisons with the zero-shot baselines should nonetheless be read
with the retrieval component in mind.
\item \textit{Extraction-bounded measurement.} A large fraction of residual
errors are answer-extraction or generation-length failures; a manual annotation
study is required to separate these from substantive reasoning errors. The
rule-based taxonomy is automatically assigned and should be read as diagnostic
rather than definitive. In particular, no human or expert annotation of the error
categories was performed; quantifying the precise share of genuine reasoning
failures would require such a study, which we leave to future work.
\item \textit{Budget-conditional reward finding.} The advantage of the
correctness-only reward holds at a 150-step budget and may not persist at longer
horizons.
\item \textit{Single in-domain benchmark.} The primary results use the single
\dataset{} train/validation split; the MMLU and FOLIO evaluations are external
sanity checks rather than a broad generalisation study.
\item \textit{Logic gap.} The physics--logic gap and the logical-fallacy error
rate indicate that standard SFT and GRPO are insufficient for formal symbolic
logic at 4B scale; step-level supervision
\citep{Lightman2023,Wang2024mathshepherd} or solver-guided training are the most
direct avenues for improvement. Because the measured gains are driven primarily
by the physics domain while logic accuracy remains below 23\%, \modelname{}
should be read as a low-cost empirical study of what curriculum and retrieval
achieve at this scale rather than as a complete solution to mixed-domain
reasoning.
\end{itemize}

%% ============================================================
\section{Conclusion}\label{sec8}
%% ============================================================

This paper presented \modelname{}, a low-cost, retrieval-augmented expert
reasoning system for mixed physics--logic question answering that combines
curriculum SFT, GRPO with lightweight reward shaping, retrieval-augmented
inference, symbolic verification, and an optional Dual-LoRA routing agent. Under a uniform single-pass protocol, the strongest
configuration reaches 47.47\% overall accuracy on the \dataset{} validation set,
a 17.51~pp improvement over the \basemodelname{} zero-shot configuration, using a
4B open model and under US\$64 of total compute. Across the three training
stages, the initial supervised fine-tuning delivers the highest accuracy gain per
training dollar (Table~\ref{tab9}); together with training-free retrieval, it
recovers most of the attainable gain for under US\$4 of fine-tuning, whereas the
later GRPO stages add compute without a statistically separable benefit. A
controlled
decomposition attributes $+9.22$~pp of the improvement to retrieval-augmented
exemplars and a further $+5.07$~pp to supervised fine-tuning, while the
GRPO-over-SFT increment is not statistically separable at the validation-set
size. A fair multi-prompt re-evaluation places the 7B-class baselines at
11.5--29.0\% and supplying them with the same retrieval does not close the gap,
external evaluations show positive transfer on MMLU physics ($+6.33$~pp) and
FOLIO logic ($+13.80$~pp over zero-shot), and a tool-augmented variant produces
only a marginal change. The error analysis indicates that answer-extraction and
generation-length failures account for a large share of residual errors and that
formal logic is the principal reasoning bottleneck. In sum, \modelname{}
substantially improves low-cost scientific reasoning---chiefly in the physics
domain---while revealing that formal logic remains the principal bottleneck for
4B-scale models, and that retrieval, rather than additional weight training,
supplies most of the attainable gain. Future work includes a
manual error-annotation study, a unified single-protocol re-evaluation that
removes the retrieval asymmetry, scaling to 7B and 14B checkpoints, cross-dataset
evaluation on SciBench \citep{Wang2024scibench} and GPQA \citep{Rein2023}, and
tighter integration of solver-guided training signals for the logic domain.

%% ---- Back matter ----
\bmsubsection*{Author Contributions}
The author designed the study, implemented the pipeline, conducted the
experiments, analysed the results, and wrote the manuscript.

\bmsubsection*{Acknowledgments}
Withheld for review.

\bmsubsection*{Funding}
The author received no specific funding for this work.

\bmsubsection*{Conflicts of Interest}
The author declares no conflicts of interest.

\bmsubsection*{Data Availability Statement}
To support reproducibility, an anonymised repository for review
(\url{https://github.com/ANONYMISED-FOR-REVIEW}) provides the training and
evaluation scripts, the trained LoRA adapters, the exact 90/10 train/validation
split, the raw per-sample single-pass predictions for every configuration, the
prompt templates and the answer extractor and verifier (including the
leakage-audit script \texttt{step10\_leakage\_check.py}), and the compute and
cost logs underlying Table~\ref{tab9}. A permanent, citable archive
(OSF or Zenodo) will be deposited on acceptance. The \dataset{} competition data
are governed by the challenge access policy and are not redistributed here; the
MMLU \citep{Hendrycks2021} and FOLIO \citep{Han2022folio} benchmarks used for
external evaluation are publicly available from their original sources.

%% NOTE: USG.cls already issues \bibliographystyle for [APA]; do NOT add another
%% (a 2nd \bibliographystyle => BibTeX "Illegal, another \bibstyle command").
\bibliography{references}

\appendix
\bmsection{Prompt Templates and Answer Extraction\label{appA}}

The baseline systems in Table~\ref{tab3} are evaluated under a fair
multi-prompt protocol with three answer formats, and the best score per model is
reported. Each prompt is supplied through the native chat template of the
corresponding model, as a system message followed by a user message containing
the question. The three system prompts are reproduced below.

\begin{lstlisting}[caption={XML prompt.},captionpos=b]
You are an expert in Logic and Physics. Think step by step and
respond in the EXACT format:
<reasoning>
[Your step-by-step reasoning]
</reasoning>
<answer>
[Final answer: letter / Yes/No/Unknown / number+unit]
</answer>
<explanation>
[Concise explanation]
</explanation>
\end{lstlisting}

\begin{lstlisting}[caption={Boxed chain-of-thought prompt.},captionpos=b]
You are an expert in Logic and Physics. Solve the problem step by
step. At the end, write your final answer inside \boxed{...}.
For MCQ, write the letter. For yes/no questions, write Yes, No, or
Unknown. For numerical problems, include units.
\end{lstlisting}

\begin{lstlisting}[caption={Plain chain-of-thought prompt.},captionpos=b]
You are an expert in Logic and Physics. Solve the problem step by
step, showing all reasoning. End your response with:
'The answer is: [your answer]'
For MCQ, the answer is just the letter (A/B/C/D). For yes/no
questions, answer Yes, No, or Unknown. For numerical problems,
include units.
\end{lstlisting}

A single extractor is applied to every response in the following priority
order: (i) the content of an \texttt{<answer>}\,\ldots\,\texttt{</answer>}
tag; (ii) the content of a \texttt{\textbackslash boxed\{\ldots\}} expression;
(iii) the text following an ``Answer:'' or ``The answer is:'' marker; and
(iv) the last non-empty line. Physics answers are compared after SI-unit
normalisation with a $5\%$ relative tolerance, and logic answers are compared by
normalised string matching, with the Z3 solver invoked when a parseable formal
representation is available.

\end{document}
